% from ega2-1-fr.tex
\oldpage[II]{162}
$h_j^{m_jk}\in T_1^kT\subset T'$. Il y a par suite un entier $m_0\geq k$ tel
que $h_j^m\in T_1^{mn_j}$ pour $m\geq m_0$. Cela étant, si $q$ est le plus
grand des entiers $n_j$, le nombre $n_0=qsm_0+k$ répond à la question. En
effet, un élément de $S_n$, pour $n\geq n_0$, est somme de monômes appartenant
à $T_1^\alpha u$, où $u$ est un produit de puissances des $h_j$~; si
$\alpha\geq k$, il résulte de ce qui précède que $T_1^\alpha u\subset T_1^n$~;
dans le cas contraire, un des exposants des $h_j$ est $\geq m_0$, donc
$u\in T_1^\beta v$, où $\beta\geq k$ et $v$ est encore un produit de
puissances des $h_j$~; on est alors ramené au cas précédent, et on voit
finalement que $T_1^\alpha u\subset T_1^n$ dans tous les cas.

\begin{remark}[8.2.15]
\label{II.8.2.15-fr}
La condition $S_n=S_1^n$ pour $n\geq n_0$ entraîne évidemment
$S_{n+1}=S_1S_n$ pour $n\geq n_0$, mais la réciproque est inexacte, même
quand on suppose $S$ noethérien. Par exemple, soient $K$ un corps,
$A=K[x]$, $B=K[y]/y^2K[y]$, où $x$ et $y$ sont deux indéterminées, $x$ étant
prise de degré $1$ et $y$ de degré $2$, et soit $S=A\otimes_KB$, de sorte que
$S$ est une algèbre graduée sur $K$ ayant une base formée des éléments $1$,
$x^n$ ($n\geq1$) et $x^ny$ ($n\geq0$). Il est immédiat que
$S_{n+1}=S_1S_n$ pour $n\geq2$, mais $S_1^n=Kx^n$ tandis que
$S_n=Kx^n+Kx^{n-2}y$ pour $n\geq2$.
\end{remark}

\subsection{Cônes projetants}
\label{subsection:II.8.3-fr}

\begin{env}[8.3.1]
\label{II.8.3.1-fr}
Soit $Y$ un préschéma~; dans tout ce numéro, il ne sera question que de
\emph{$Y$-préschémas} et de \emph{$Y$-morphismes}. Soit $\mathcal{S}$ une
$\mathcal{O}_Y$-Algèbre quasi-cohérente graduée à \emph{degrés positifs}~;
\emph{nous supposerons en outre que l'on a $\mathcal{S}_0=\mathcal{O}_Y$}.
Conformément aux notations introduites en
(\hyperref[II.8.2.2-fr]{8.2.2}), nous poserons
\begin{equation}
\label{II.8.3.1.1-fr}
  \widehat{\mathcal{S}}=\mathcal{S}[z]
  =\mathcal{S}\otimes_{\mathcal{O}_Y}\mathcal{O}_Y[z]
\tag{8.3.1.1}
\end{equation}
que l'on considère comme $\mathcal{O}_Y$-Algèbre graduée à degrés positifs en
définissant les degrés par la formule
(\hyperref[II.8.2.2.2-fr]{8.2.2.2}), de sorte que pour tout ouvert affine $U$
de $Y$, on a
\[
  \Gamma(U,\widehat{\mathcal{S}})=(\Gamma(U,\mathcal{S}))[z].
\]
Dans ce qui suit, nous poserons
\begin{equation}
\label{II.8.3.1.2-fr}
  X=\operatorname{Proj}(\mathcal{S}),\qquad
  C=\operatorname{Spec}(\mathcal{S}),\qquad
  \widehat{C}=\operatorname{Proj}(\widehat{\mathcal{S}})
\tag{8.3.1.2}
\end{equation}
(où dans la définition de $C$, $\mathcal{S}$ est considérée comme
$\mathcal{O}_Y$-Algèbre non graduée), et nous dirons que $C$ (resp.
$\widehat{C}$) est le \emph{cône affine} (resp. \emph{cône projectif}) défini
par $\mathcal{S}$~; on dira aussi «~cône~» au lieu de «~cône affine~». Par
abus de langage, nous dirons encore que $C$ (resp. $\widehat{C}$) est le
\emph{cône projetant affine de $X$} (resp. le \emph{cône projetant projectif
de $X$}), étant sous-entendu que le préschéma $X$ est donné sous la forme
$\operatorname{Proj}(\mathcal{S})$~; enfin, nous dirons que $\widehat{C}$ est
la \emph{fermeture projective} de $C$ (étant entendu que la donnée de
$\mathcal{S}$ figure dans la structure de $C$).
\end{env}

\begin{proposition}[8.3.2]
\label{II.8.3.2-fr}
Il existe des $Y$-morphismes canoniques
\begin{equation}
\label{II.8.3.2.1-fr}
  Y\xrightarrow{\varepsilon}C\xrightarrow{i}\widehat{C}
\tag{8.3.2.1}
\end{equation}
\begin{equation}
\label{II.8.3.2.2-fr}
  X\xrightarrow{j}\widehat{C}
\tag{8.3.2.2}
\end{equation}

